\documentclass[a4paper]{article}
\usepackage{bookmark}
\usepackage[italian]{babel}
\usepackage{geometry}
\usepackage{hyperref}
\usepackage{sectsty}
\usepackage{parskip}

\hypersetup{hidelinks}
\allsectionsfont{\centering}

\begin{document}
    
    \title{Analisi dell'amplificatore Miller}
    \author{Luca Brescia}
    \date{\today}
    \maketitle

    \section*{Obiettivo}
        Creazione di uno standard per le simulazioni dei circuiti della tesi. \\
        Utilizzare prima il circuito di un amplificatore Miller, in configurazione di inseguitore (buffer), a cui vanno effettuate le seguenti simulazioni:

        \begin{itemize}
            \item Analis DC: in cui bisogna osservare l'assestamento dell'uscita al valore prefissato di ingresso.
            \item Analisi AC: fare sweep in frequenza da 1 MHz a 1-2-4 GHz
        \end{itemize}

        L'ingresso dovrà essere impostato su 0.5Vpp e ampiezza 1Vpp basato su una tensione di alimentazione di Vdd=1V.

        Inoltre si dovranno effettuare le simulazioni con le seguenti configurazioni:

        \begin{itemize}
            \item 1 stadio;
            \item 2 stadi;
            \item Eff. corner;
            \item PVT
            \item Montecarlo;
        \end{itemize}

    \section*{Svolgimento}
        \subsection*{Miller a singolo stadio}
            Si è partiti con la simulazione di un amplificatore Miller a singolo stadio. La prima analisi che si è effettuata è l'andamento dell'uscita a fronte di un ingresso costante a $V_+ = V_- = \frac{Vdd}{2}$ per dimensionare i transistori MOS di modo da ottenere $V_o \simeq \frac{Vdd}{2}$. \\
            Si è scelto di dimensionare i transistori con una L non minima, poiché l'applicazione è analogica, e si è deciso di partire con una $L = 4\mu m$. Con tale valore si è effettuato uno sweep in W dei transistori NMOS, detta \textit{wn\_miller} e, successivamente, si è effettuato uno sweep in W dei transistori PMOS, detta \textit{wp\_miller}.
            I valori di sweep utilizzati sono riportati tabella \ref{tab:sweep_wn_wp}:
            \begin{table}[]
                \centering
                \begin{tabular}{c|c|c|}
                    \cline{2-3}
                                                     & da        & a          \\ \hline
                    \multicolumn{1}{|c|}{wn\_miller} & $8\mu m$  & $300\mu m$ \\ \hline
                    \multicolumn{1}{|c|}{wp\_miller} & $16\mu m$ & $300\mu m$ \\ \hline
                    \end{tabular}
                \label{tab:sweep_wn_wp}
                \caption[short]{Valori di sweep per le variabili di W dei transistori MOS per l'amplificatore Miller a singolo stadio}
            \end{table}
            Si sono scelti i valori $wn\_miller = 24\mu m$ e $wp_miller = 72\mu m$. 
            Con tali geometrie si ottiene $V_o \simeq \frac{Vdd}{2}$ con tutti i transistori in regione 2 (saturazione) tranne per il generatore di corrente, polarizzato con una tensione $V_b = 800mV$, che resta in regione 1 (lineare).
            Per la scelta dell tensione $V_b$ di polarizzazione del generatore di corrente, transistore M3, si è effettuato uno sweep in tensione da $V_{th} \simeq 457mV$ a $Vdd$. Si è osservato un mantenimento della regione 1 per tutto lo sweep di tensioni soprasoglia e che, come ci si aspetterebbe, la tensione di uscita cala all'aumentare della tensione $V_b$.


    \section*{Appunti incontro Richelli}
    Il dimensionamento dei componenti devo farlo in modo tale da fissare il dc operating point in SAT per tutti i MOS.
    Quindi prima dimensiono in modo da avere SAT anche per il gen di corrente. Una volta che so che nel punto di lavoro il MOS è in SAT posso fare sweep delle tensioni V+ e V- per vedere la caratteristica.
    Quello che mi aspetto è che la Vd del gen di corrente dipende da Vs della coppia differenziale e quindi da V+ e V-. Per cui sicuramente andrà in LIN il gen di corrente, tuttavia devo dimensionare in modo da avere ro grande e gm piccola e questo dipende dalle geometrie (da vedere bene come sono legati a L e W).

    Inoltre, per il dimensionamento del NAUTA, sono tutti inverter per cui devo mantenere le caratteristiche dell'inverter ma siamo nel mondo analogico per cui vanno dimensionati ragionevolmente grandi. 
    Da tenere in considerazione le geometrie del miller per cui gli NMOS sono dimensionati più o meno uguali. I PMOS devono essere più grandi (sempre come nel miller) quindi fare riferimento alle geometrie del paper EMCCompo.
    Controllare la caratteristica degli inverter base che deve essere rispecchiata anche per questa tipologia di inverter. 
            
\end{document}